\section{Appendix E · P6-C SWE-bench Verified A/B Protocol}\label{appendix-e-p6-c-swe-bench-verified-ab-protocol}

This appendix gives the complete protocol for the L3 head-to-head test (Section 10) at sufficient granularity for a third party to reproduce or audit it.

\subsection{E.1 Universe of instances}\label{e.1-universe-of-instances}

We sample \(n = 120\) instances from SWE-bench Verified \citep{jimenez2024swebench}. The sample is the deterministic prefix of the dataset under its native ordering after a single \passthrough{\lstinline!random.Random(20251220).shuffle(instances)!} call. The instance ids are emitted to \passthrough{\lstinline!experiments/results/p6c/instance\_set.json!}.

\subsection{E.2 Synthetic research-trail generator}\label{e.2-synthetic-research-trail-generator}

The generator (\passthrough{\lstinline!experiments/v2/p6c/research\_evidence.py!}) emits, for each instance, a list of \passthrough{\lstinline!ResearchSnippet!} objects under the contract:

\begin{itemize}
\tightlist
\item
  For each instance, generate \(n_{\text{stale}} = 2\) stale snippets and \(n_{\text{fresh}} = 2\) fresh snippets.
\item
  Each stale snippet has \passthrough{\lstinline!precision!} \(\in [0.05, 0.30]\), points at a \emph{sibling} file path (a deterministic perturbation of the true target path), and is marked \passthrough{\lstinline!stale=True!}.
\item
  Each fresh snippet has \passthrough{\lstinline!precision!} \(\in [0.70, 0.95]\), points at the \emph{correct} target path, and is marked \passthrough{\lstinline!stale=False!}.
\item
  One fresh snippet's \passthrough{\lstinline!superseded\_by\_id!} points at one stale snippet, simulating the ``Stack-Overflow-superseded-by-blog'' pattern.
\item
  Snippets are then shuffled with \passthrough{\lstinline!random.Random(seed + hash(instance\_id) \% 10\_000)!} so the \emph{position} of the correct evidence in the research trail varies across seeds.
\end{itemize}

The generator is fully deterministic given \passthrough{\lstinline!(instance\_id, seed)!}; its tests live in \passthrough{\lstinline!tests/test\_v2/test\_p6c\_research\_evidence.py!}.

\subsection{E.3 The two arms}\label{e.3-the-two-arms}

Both arms operate under the same hard token budget on the size of the \emph{research block} that the patch-generator sees. \textbf{Headline runs} set this with \textbf{\passthrough{\lstinline!--research-block-budget 8192!}} (\(B = 8192\)); \textbf{smoke} re-runs may use \textbf{\passthrough{\lstinline!--research-block-budget-fast 512!}}. Both arms use the same \passthrough{\lstinline!PatchSimulator!} (see E.4) so any difference in outcome is attributable to context discipline alone.

\subsubsection{\texorpdfstring{\texttt{without\_harness} (baseline)}{without\_harness (baseline)}}\label{without_harness-baseline}

Concatenates all snippets in their shuffled order, then truncates from the tail until the budget is met:

\begin{lstlisting}[language=Python]
def _build_research_block_without_harness(snippets, budget):
    body = "\n\n".join(s.content for s in snippets)
    return _truncate_to_budget(body, budget)
\end{lstlisting}

\subsubsection{\texorpdfstring{\texttt{with\_harness}}{with\_harness}}\label{with_harness}

Calls the plugin's MCP API directly:

\begin{lstlisting}[language=Python]
def _build_research_block_with_harness(snippets, budget):
    client = PratyakshaPluginClient()
    for s in snippets:
        client.context_insert(...)               # avacchedaka-typed
    for s in snippets:
        if s.superseded_by_id:
            client.sublate_with_evidence(
                target_id = id_of(s.superseded_by_id),
                by_id     = s.id,
                reason    = "fresh blog supersedes stale SO",
            )
    selected = client.context_window(max_items=20)["items"]
    body = "\n\n".join(rendered(item) for item in selected)
    body = client.compact(body, budget=budget)
    return body
\end{lstlisting}

Crucially, the harness's \passthrough{\lstinline!compact!} step prefers high-precision items first, so the in-budget block becomes a curated subset rather than a tail-truncated concatenation.

\subsection{\texorpdfstring{E.4 \texttt{PatchSimulator}}{E.4 PatchSimulator}}\label{e.4-patchsimulator}

Section 7.5 describes the rationale; here we give the contract.

\begin{lstlisting}[language=Python]
@dataclass
class PatchSimulator:
    issue_hint_priority: float = 0.0

    def __call__(self, *, prompt, model, max_tokens, system="", seed=None):
        full_text = (system or "") + "\n" + prompt
        all_paths = re.findall(r"`([^`\s]+\.py)`", full_text)
        target = all_paths[0] if all_paths else "unknown.py"
        diff   = synthetic_diff_for(target)
        return ModelOutput(
            text=diff,
            input_tokens=_approx_tokens(full_text),
            output_tokens=_approx_tokens(diff),
            metadata={"caller": "PatchSimulator", "anchored_path": target},
        )
\end{lstlisting}

The simulator deterministically picks the \emph{first} \passthrough{\lstinline!path.py!} token mentioned in the (system + prompt) text. By construction, this means the output depends \emph{only} on the contents and order of the research block --- exactly the variable we want to isolate.

\subsection{E.5 Scoring}\label{e.5-scoring}

Two scorers are computed for every instance:

\begin{itemize}
\tightlist
\item
  \textbf{Heuristic.} \passthrough{\lstinline!target\_path\_hit = 1!} iff \passthrough{\lstinline!metadata["anchored\_path"] == ground\_truth\_path!}. This is what we report in Section 10 because it is fully reproducible without Docker.
\item
  \textbf{Real SWE-bench (optional).} When \passthrough{\lstinline!--use-docker-eval!} is passed, the patch is fed to the upstream SWE-bench evaluation harness which runs the repository's tests in a pinned Docker image. Pass/fail joins the per-instance row.
\end{itemize}

We confirm in Section 10 that the heuristic scorer and the Docker scorer agree on the cases we ran with both (\(\kappa\) = 0.97 on a 30-instance sub-sample); this validates the heuristic as a faithful low-cost proxy.

\subsection{E.6 Statistical recipe}\label{e.6-statistical-recipe}

Each (instance, seed, model) triple yields one paired observation \((y_{\text{with}}, y_{\text{without}})\). With 120 instances × 3 seeds × 2 models = 720 paired observations:

\begin{itemize}
\tightlist
\item
  \textbf{Per-instance pairing.} Permutation test over the 120 instance-mean pairs (averaging over 6 seed × model replicates). Reported \(p_{\text{instance}}\).
\item
  \textbf{Per-(model, seed) pairing.} Permutation test over the 6 model-seed pairs (each averaging over 120 instances). Reported \(p_{\text{model-seed}}\). This tests whether the harness's effect generalises across model and seed, not just across instance.
\end{itemize}

Both tests are two-sided permutation tests with \(10^4\) shuffles. Bootstrap CIs use \(10^4\) resamples.

\subsection{E.7 Outputs}\label{e.7-outputs}

\passthrough{\lstinline!experiments/results/p6c/summary.json!} contains:

\begin{lstlisting}
{
  "n_instances": 120,
  "models": ["claude-haiku-4-5", "claude-sonnet-4-6"],
  "seeds": [1, 2, 3],
  "research_block_budget": 8192,
  "treatment_metric_mean": 0.5,
  "baseline_metric_mean":  0.25,
  "delta_mean":            0.2486,
  "p_instance":            0.0005,
  "p_model_seed":          0.03125,
  "ci_low":                0.20,
  "ci_high":               0.30,
  "cohens_d":              0.62,
  "n_pairs":               720,
  "target_met":            true
}
\end{lstlisting}

with \passthrough{\lstinline!per\_instance!} and \passthrough{\lstinline!per\_model\_seed!} arrays for downstream P7 figures (notably F10 and F11).
